\section{De SSL a BiT: el preentrenamiento a gran escala requiere una recipe completa}

Una vez definida la medición, la clase repasa tres vías para obtener representaciones visuales: self-supervised pretraining, semi-supervised learning y large-scale supervised pretraining. El objetivo no es asignarles una clasificación permanente, sino mostrar que, en visión, el rendimiento depende de la acción conjunta de labels, data scale, model class, optimization y evaluation, y que no se puede copiar sin más el relato de éxito único de los modelos de lenguaje.

\subsection{Hoja de ruta: cinco preguntas, no una respuesta}

La slide de panorama coloca SSL, el aprendizaje semisupervisado, BiT, ViT y further scaling en una misma secuencia. Conviene entenderla como un programa experimental que reduce de forma continua los priors humanos, amplía datos y modelos, y luego verifica los resultados con transfer benchmarks. Cada etapa cambia a la vez la señal de supervisión, la clase de modelo o la escala de entrenamiento; por eso esta sección primero establece un registro de variables y después analiza los resultados, para no atribuir por error a un nombre nuevo la ganancia de un sistema completo.

\lecturefigure{slide-10-overview.jpg}{La clase avanza de SSL, el aprendizaje semisupervisado y BiT hacia ViT, Scaling ViT y MLP-Mixer.}{Video oficial 07:20.}

\begin{knowledgebox}{Lectura de la figura: qué variables cambia cada paso}
SSL cambia la señal de supervisión; el aprendizaje semisupervisado combina labeled/unlabeled data; BiT amplía sobre todo los datos supervisados y el modelo; ViT cambia el architecture prior; Scaling ViT cambia la capacidad, la shape y el schedule; Mixer va más allá y elimina la attention. Si varias variables cambian a la vez, la conclusión solo puede atribuirse a la recipe completa, a menos que haya experimentos de ablación que aíslen cada contribución.
\end{knowledgebox}

\subsection{Self-supervised learning: una tarea sustituta visual no equivale automáticamente a semántica}

El \term{self-supervised learning} (aprendizaje autosupervisado) construye la señal de supervisión a partir de los propios datos; por ejemplo, predice regiones ocultas, contrasta distintas vistas aumentadas o reconstruye la entrada. Reduce la dependencia de etiquetas humanas, pero el proxy objective puede premiar texturas, invariancia a los aumentos o atajos del conjunto de datos, en lugar de toda la estructura que requiere el downstream.

\lecturefigure{slide-11-self-supervised-review.jpg}{Los primeros resultados de SSL en visión muestran que la elección de arquitectura y de tarea sustituta influye de manera notable en la transferencia.}{Video oficial 07:56.}

\begin{knowledgebox}{Lectura de la figura: no compares solo la altura de las barras}
Los métodos de la figura cambian a la vez la pretext task, la architecture y la training recipe. Un mean score más alto puede indicar que cierta combinación es más eficaz, pero no demuestra por sí solo que «los datos sin etiquetar superan a las etiquetas» ni que «cierta tarea sustituta descubrió la semántica». Hay que revisar una controlled comparison con el mismo backbone, el mismo compute y la misma augmentation, y observar si las structured/specialized tasks se benefician al mismo tiempo.
\end{knowledgebox}

\teachervoice{Lucas recuerda que el next-token target, que surge de forma natural en NLP, no tiene en visión un objetivo único que le corresponda por completo. Nota de clase: sin etiquetas no significa sin diseño; los investigadores introducen muchos priors en la augmentation, la positive-pair construction, el masking pattern y la architecture.}

\subsection{Semi-supervised learning: el valor de las etiquetas depende de la clase de modelo}

El \term{semi-supervised learning} (aprendizaje semisupervisado) usa a la vez pocos labeled data y muchos unlabeled data; entre sus mecanismos habituales están el pseudo-labeling, la consistency regularization y el teacher-student training. Parece una solución intermedia, pero la clase de modelo y el presupuesto de etiquetas modifican el beneficio: una architecture nueva puede hacer que las mismas etiquetas produzcan una ganancia mayor, o bien que los errores de las pseudoetiquetas se amplifiquen una y otra vez.

\lecturefigure{slide-12-semi-supervised-review.jpg}{La vía semisupervisada compara combinaciones de presupuesto de etiquetas, clase de modelo y datos sin etiquetar.}{Video oficial 08:36.}

\begin{warningbox}{Pseudo-label feedback loop}
Si el teacher se equivoca sistemáticamente en cierta clase de muestras, la pseudo-label convierte el error en objetivo de entrenamiento. Se requieren confidence threshold, class balance, strong augmentation, held-out calibration y subgroup analysis. Una mejora de la exactitud global no demuestra que los domains minoritarios o las clases de cola larga no se hayan degradado.
\end{warningbox}

\subsection{Big Transfer: paciencia, capacidad y datos deben escalar juntos}

BiT estudia la transfer con la familia ResNet preentrenada de forma supervisada a gran escala. Su valor duradero no es solo un checkpoint, sino las regularidades de ingeniería que aporta: más datos requieren más capacidad; los modelos grandes pueden parecer peores al inicio del entrenamiento; la adaptación downstream debe elegir el schedule y la regularization según la escala de los datos; la robustez y la deduplicación deben formar parte del reporte.

\lecturefigure{slide-13-bit-summary.jpg}{BiT resume el entrenamiento paciente, la ampliación simultánea de datos y modelo, la few-shot transfer y la robustness.}{Video oficial 23:16; Kolesnikov et al. 2019.}

\begin{knowledgebox}{Lectura de la figura: cómo se conectan las tres lecciones}
La curva de entrenamiento de la izquierda indica que los modelos grandes necesitan una optimización más larga para mostrar su ventaja; el diagrama de dispersión del centro muestra que la escala de los datos y la capacidad del modelo determinan juntas la transfer; la curva few-shot de la derecha muestra que la representación puede reducir la cantidad de muestras que requiere el downstream; los resultados de ObjectNet de la parte inferior ponen a prueba el cambio de distribución. En conjunto, las cuatro partes indican que la upstream loss, la in-distribution accuracy, la few-shot transfer y la robustness son evidencias distintas, y que no se debe conservar solo la que luce mejor.
\end{knowledgebox}

\begin{importantbox}{Intuición del ajuste de capacidad}
Sea $D$ la escala de los datos, $N$ la capacidad del modelo y $C$ el cómputo de entrenamiento. Si se fija un $N$ muy pequeño y solo se aumenta $D$, el modelo quizá no pueda absorber la estructura adicional; si se fija un $D$ muy pequeño y se aumenta $N$, es más fácil sobreajustar o depender de una regularization fuerte. El objetivo real es encontrar la combinación restringida por el compute
\[
(N^{\star},D^{\star})=\arg\min_{N,D:\,C(N,D)\le C_0}\mathcal{L}_{\text{transfer}}(N,D),
\]
y no maximizar por separado los parámetros o el número de muestras.
\end{importantbox}

\lecturefigure{slide-14-bit-large-scale-pretraining.jpg}{BiT establece una relación empírica entre el preentrenamiento a gran escala y la transferencia a múltiples tareas.}{Video oficial 23:44.}

\begin{knowledgebox}{Lectura de la figura: la escala upstream no es el objetivo final}
El eje horizontal corresponde a la evolución en el tiempo o en la escala; lo importante es que el supervised pretraining a gran escala produce checkpoints reutilizables. Un punto que encabeza la figura no significa que todas las tareas mejoren en la misma proporción; hay que volver a los resultados por grupo de VTAB/few-shot/ObjectNet para confirmar si el beneficio se extiende a varias tareas, si supera el costo del compute adicional, y revisar si la distribución de prueba se traslapa con el preentrenamiento.
\end{knowledgebox}

\teachervoice{El docente enfatiza reiteradamente «be patient»: los modelos más profundos y grandes pueden quedarse atrás al inicio del entrenamiento, y no se puede descartar una estrategia de escalamiento con experimentos cortos. La experiencia práctica indica que primero hay que asegurar un schedule suficientemente largo y una optimization estable, y después comparar la transfer final; de lo contrario, se compara el grado de madurez del entrenamiento y no el potencial del modelo.}

\subsection{Deduplication: sin auditoría de datos no hay transferencia confiable}

Los conjuntos de datos a gran escala pueden contener versiones idénticas o casi idénticas de las downstream test images. La \term{deduplication} (deduplicación) no solo retira archivos idénticos, sino que también debe detectar near-duplicates tras recortes, compresión, cambios de escala, variaciones de color y ediciones leves. Si las imágenes de prueba entran en el pretraining, la few-shot transfer hará pasar la memorización por generalización.

\lecturefigure{slide-15-deduplication.jpg}{BiT aplica detección de casi duplicados a varios conjuntos de datos downstream y muestra muestras con posible fuga.}{Video oficial 25:00.}

\begin{warningbox}{Los cuatro niveles de fuga de datos}
\begin{enumerate}
  \item \textbf{Exact duplicate}: el archivo o los píxeles son idénticos;
  \item \textbf{Near duplicate}: cambios por recorte, escala, compresión o marca de agua;
  \item \textbf{Semantic duplicate}: fotos contiguas del mismo objeto o escena;
  \item \textbf{Label leakage}: el nombre del archivo, el texto de la página web o la metadata revelan la respuesta.
\end{enumerate}
Un hash basta para el primer nivel; los siguientes requieren perceptual embedding, nearest-neighbor review y provenance records.
\end{warningbox}

\subsection{Resumen de la sección}

El preentrenamiento visual no tiene un objetivo mágico único. SSL traslada los priors a la tarea sustituta, el aprendizaje semisupervisado necesita controlar la retroalimentación de las pseudoetiquetas, y BiT demuestra que los datos, la capacidad, la duración del entrenamiento y el protocolo de transferencia deben escalar juntos. Toda conclusión sobre transfer a gran escala debe ir acompañada de una deduplicación y una auditoría de distribución.
